\documentclass[11pt,a4paper]{article}

\usepackage{amsmath,amssymb,amsthm}
\usepackage{mathtools}
\usepackage{hyperref}
\usepackage{geometry}
\geometry{margin=2.5cm}

\newtheorem{theorem}{Theorem}[section]
\newtheorem{proposition}[theorem]{Proposition}
\newtheorem{lemma}[theorem]{Lemma}
\newtheorem{corollary}[theorem]{Corollary}
\newtheorem{definition}[theorem]{Definition}
\theoremstyle{remark}
\newtheorem{remark}[theorem]{Remark}

\DeclareMathOperator{\Sym}{Sym}
\DeclareMathOperator{\ad}{ad}
\newcommand{\heff}{\mathcal{H}_{\mathrm{eff}}}
\newcommand{\Leff}{L_{\mathrm{eff}}}
\newcommand{\LWeil}{L_{\mathrm{Weil}}}
\newcommand{\gsp}{g_{\mathrm{sp}}}
\newcommand{\Vr}{V_{\rho}}
\newcommand{\Wsp}{W_{\mathrm{sp}}}
\newcommand{\suTwo}{\mathfrak{su}(2)}
\newcommand{\heisR}{\mathrm{Heis}_{3}(\mathbb{R})}
\newcommand{\heisZ}{\mathrm{Heis}_{3}(\mathbb{Z}/q\mathbb{Z})}
\newcommand{\Csu}{C_{\suTwo}}
\newcommand{\Az}{A_{Z}}
\newcommand{\AH}{A_{H}}
\newcommand{\tX}{\widetilde{X}}
\newcommand{\tY}{\widetilde{Y}}
\newcommand{\tZ}{\widetilde{Z}}
\newcommand{\sigmatwo}{\sigma_{2}}

\bibliographystyle{plain}
\usepackage{doi}

\begin{document}

  \title{Sub-Principal Symbol of the Effective Operator and Casimir Rigidity of the Central Direction\\
  \large Q8 of the Cosmochrony Spectral Geometry Programme}
  \author{J\'er\^ome Beau\\
  \small Independent Researcher\\
  \small \texttt{jerome.beau@cosmochrony.org}}
  \date{2026}
  \maketitle

  \begin{abstract}
    Paper Q7~\cite{Beau2026q7} reduced the open problem Q5b-O2 of~\cite{Beau2026q5b} to a single
    computable criterion: whether the sub-principal symbol of the effective operator $\Leff$ in the
    central direction $\tZ = [\tX, \tY]$ of $\heisR$ carries coefficient $\Az = \Csu = 2$, the eigenvalue
    of the $\suTwo$-Casimir on the spin-$1$ module $\Sym^{2}(\Vr)$.
    The no-cross-terms part of that criterion is already a theorem (Q7 Proposition~6.1).
    The present paper settles the isotropy condition $\AH = \Az$ by a structural argument that
    does not require a full computation of the hypoelliptic symbol.
    The key observation is that the sub-principal term in $\Leff$ along $\tZ$ has a unique algebraic
    source: the Heisenberg commutator $[\tX, \tY] = \tZ$, which reflects the nilpotency class~$2$
    of $\heisR$ and cannot be mimicked by any other generator.
    Under the equivariant bridge $\varphi \colon \Sym^{2}(\Vr) \xrightarrow{\;\sim\;} \Wsp$ established
    in Q7 (unique up to positive scalar by Schur's lemma), the image of this commutator term under
    $\varphi$ is constrained to be proportional to the unique $\suTwo$-invariant quadratic form on
    $\Sym^{2}(\Vr)$, namely the Casimir $\Csu = 2 \cdot \mathrm{Id}$.
    There is no free scalar: the normalisation is fixed by the same Casimir that sets $\Az = 2$ in
    Q7 Remark~6.3.
    The result is therefore: \emph{given bridge non-obstruction} (proved in Q9~\cite{Beau2026q9}),
    $\Az = \Csu = 2$ unconditionally under the Q5a hypotheses.
    The lifting hypothesis~[H-lift], previously an additional assumption of Q5b, is now a theorem
    of Q9~\cite{Beau2026q9} under the Q5a hypotheses alone, so no extra condition is needed.
    The effective spatial co-metric is $\gsp = \AH(k_X^{2} + k_Y^{2}) + 2\,k_Z^{2}$ with $\AH \to 2$
    as $q \to \infty$, proved in Q10~\cite{Beau2026q10} and U1~\cite{Beau2026u1} under the Q5a
    hypotheses and the O-series spectral universality~[U].
    Together with Q5b Theorem~6.1, this yields a non-degenerate rank-$4$ Lorentzian metric on
    $\mathbb{R}_{\tau} \times \heisR$ and closes Q5b-O2 under the Q5a hypotheses.
  \end{abstract}

  \textbf{Keywords:} Heisenberg group; Weil representation; sub-principal symbol; Casimir operator;
  $\suTwo$-module; nilpotent Lie algebra; effective metric; Cosmochrony.

  \medskip
  \noindent\textbf{MSC 2020.} 22E25, 81R05, 53B30, 35H10.

  \clearpage
  \tableofcontents
  \clearpage

  \section{Introduction}

    \subsection{Context and the reduction of Q7}

      Paper Q5b~\cite{Beau2026q5b} derives the effective Lorentzian geometry of the Cosmochrony
      framework by extracting a metric tensor from the principal symbol of the effective operator
      $\Leff$ on $\mathbb{R}_{\tau} \times \heisR$.
      The principal symbol is degenerate: the central generator $\tZ = [\tX, \tY]$ of $\heisR$ does not
      appear at leading sub-Riemannian order, entering only through sub-principal corrections.
      Open problem Q5b-O2 asks for the value of the coefficient $\Az$ of the resulting $k_Z^{2}$-term,
      which would complete the $4 \times 4$ effective metric to full rank.

      Paper Q7~\cite{Beau2026q7} established three structural results constraining the bridge
      $\varphi \colon \heff \simeq \Sym^{2}(\Vr) \to \Wsp$:
      (i)~$\suTwo \not\simeq \heisR$ as Lie algebras, so the bridge operates at the representation
      level;
      (ii)~$\varphi$ is unique up to positive scalar (Rigidity Lemma, Q7 Lemma~4.3);
      (iii)~the no-cross-terms part of Criterion~5.1 of Q7 is a theorem (Q7 Proposition~6.1).
      The remaining open part after Q7 was the isotropy condition $\AH = \Az$ and the identification
      $\Az = \Csu = 2$, now settled by Q8\textendash Q10 and U1.

    \subsection{Main result}

      \begin{theorem}[Casimir rigidity of the central direction]\label{thm:main}
        Assume the hypotheses \emph{[H1]}, \emph{[H-w]}, \emph{[H-E1]}, \emph{[C]} of\/
        $Q5a$~\cite{Beau2026q5a} and the bridge non-obstruction of\/ Q7\textendash Q9
        (proved in Q9~\cite{Beau2026q9}: the pathological Case~3 of\/ Q7 Remark~6.8 is
        excluded by the coercivity of the limit form).
        The lifting hypothesis~\emph{[H-lift]} of\/ Q5b is a theorem under these hypotheses
        \textup{(}Q9~\cite{Beau2026q9}\textup{)}, so no additional assumption is required.
        Then the sub-principal symbol of\/ $\Leff$ in the central direction satisfies
        \begin{equation}
          \sigmatwo(\Leff)\big|_{\tZ\text{-sector}} = \Az\, k_Z^{2},
          \qquad \Az = \Csu = 2.
        \end{equation}
        Consequently the effective co-metric is $g^{\mu\nu} = \mathrm{diag}(-A_{\tau}, 2, 2, 2)$
        with $\AH = \Az = 2$ \textup{(}proved in Q10~\cite{Beau2026q10} and
        U1~\cite{Beau2026u1}\textup{)}, non-degenerate of rank~$4$ with Lorentzian signature
        $(-,+,+,+)$.
        Open problem $Q5b$-$O2$ is resolved under the Q5a hypotheses and~\emph{[U]}.
      \end{theorem}

      The value $\Az = 2$ is \emph{unconditional} given bridge non-obstruction: it follows from the
      algebraic identity $\Csu|_{\Sym^{2}(\Vr)} = 2 \cdot \mathrm{Id}$ and the rigidity of $\varphi$.
      The lifting hypothesis~[H-lift] is now a theorem (Q9~\cite{Beau2026q9}), so the only remaining
      conditionality is the bridge non-obstruction (proved) and the Q5a hypotheses.

    \subsection{Organisation}

      Section~\ref{sec:setup} recalls the relevant results of Q5b and Q7 and fixes notation.
      Section~\ref{sec:source} identifies the unique algebraic source of the sub-principal $\tZ$-term
      from the nilpotency structure of $\heisR$.
      Section~\ref{sec:casimir} carries out the Casimir identification and proves Theorem~\ref{thm:main}.
      Section~\ref{sec:beals} places the result in the hypoelliptic operator framework (Beals\textendash Greiner
      calculus) as a consistency check and notes that no contributions beyond those identified in
      Section~\ref{sec:source} arise at the relevant nilpotency order.
      Section~\ref{sec:consequences} draws the consequences for Q5b-O2 and the effective metric.

  \section{Setup and Inherited Results}\label{sec:setup}

    \subsection{The effective operator and its principal symbol}

      Under the Q5a hypotheses ([H1], [H-w], [H-E1], [C]), which imply [H-lift]
      by Q9~\cite{Beau2026q9}, Q5b Theorem~5.2~\cite{Beau2026q5b} gives the effective
      operator $\Leff$ on $\mathbb{R}_{\tau} \times \heisR$ with leading-order principal symbol
      \begin{equation}
        \sigmatwo(\Leff)(p, k) = -A_{\tau}\, k_{\tau}^{2} + \AH(k_X^{2} + k_Y^{2}) + 0 \cdot k_Z^{2}
        + (\text{sub-principal in } k_Z),
      \end{equation}
      where the $k_Z^{2}$-slot is empty at principal order because the sub-Riemannian Laplacian
      $\Delta_{H} = \tX^{2} + \tY^{2}$ on $\heisR$ is degenerate in the central direction.

    \subsection{Inherited results from Q7}

      The following results from Q7~\cite{Beau2026q7} are used without re-proof:
      \begin{itemize}
        \item \emph{Symmetric square identification} (Q7 Proposition~4.1):
        $\heff \simeq \Sym^{2}(\Vr)$ as an $\suTwo$-module (spin-$1$ representation).
        \item \emph{Rigidity Lemma} (Q7 Lemma~4.3): any $\suTwo$-equivariant isomorphism
        $\varphi \colon \Sym^{2}(\Vr) \xrightarrow{\sim} \Wsp$ is unique up to a positive scalar.
        \item \emph{No-cross-terms theorem} (Q7 Proposition~6.1 and Corollary~6.2): the cross-term
        components $k_X k_Z$ and $k_Y k_Z$ of $\sigmatwo(\Leff)|_{\Sym^{2}(\Vr)}$ vanish structurally
        for all $(q, c)$.
        \item \emph{Grading correspondence} (Q7 Remark~4.2): the spin-weight decomposition
        $(\pm 1, 0)$ on $\Sym^{2}(\Vr)$ matches the Carnot degree decomposition $(1, 2)$
        of $\heisR$, placing $\tZ$ at spin-weight~$0$ and Carnot degree~$2$.
      \end{itemize}

      The open question after Q7 is therefore a single scalar: whether $\Az = \Csu = 2$.

  \section{Unique Algebraic Source of the Sub-Principal $\tZ$-Term}\label{sec:source}

    \subsection{Nilpotency and the generating role of the commutator}

      The Lie algebra $\heisR$ is nilpotent of class~$2$: the derived ideal
      $[\heisR, \heisR] = \mathbb{R}\tZ$ is one-dimensional, central, and exhausts the lower
      central series at step~$2$.
      This has a direct consequence for the symbol expansion of $\Leff$.

      \begin{proposition}[Uniqueness of the sub-principal $\tZ$-source]\label{prop:source}
        The $k_Z^{2}$-term in the symbol of\/ $\Leff$ at sub-principal order has a unique algebraic
        origin: the Heisenberg commutator $[\tX, \tY] = \tZ$.
        No other combination of generators of\/ $\heisR$ can contribute a term of the form $f(k_Z)\,k_Z^{2}$
        at this order, because $\tZ$ is central and appears in no other commutator relation.
      \end{proposition}

      \begin{proof}
        The generators $\{\tX, \tY, \tZ\}$ satisfy the sole non-trivial relation $[\tX, \tY] = \tZ$,
        with $\tZ$ central: $[\tZ, \tX] = [\tZ, \tY] = 0$.
        In the symbol calculus of a left-invariant differential operator on a graded nilpotent group,
        the sub-principal symbol at degree-$2$ Carnot order receives contributions only from the
        homogeneous degree-$2$ part of the operator and from commutator terms of degree-$(1,1)$
        generators.
        The only such commutator is $[\tX, \tY] = \tZ$.
        Since $\tZ$ is central, there is no further commutator that could generate additional $k_Z$-terms.
        The class-$2$ nilpotency is therefore the structural reason why the $k_Z^{2}$ coefficient is
        determined by a single algebraic datum.
      \end{proof}

      \begin{remark}[Contrast with higher-step nilpotent groups]
For a nilpotent group of class~$3$ or higher, additional commutators $[\tZ, \cdot]$ would be
non-trivial and could contribute further sub-principal terms.
The class-$2$ nilpotency of $\heisR$ closes the source list at a single term, making
the identification problem tractable at this order.
      \end{remark}

  \section{Casimir Identification and Proof of the Main Theorem}\label{sec:casimir}

    \subsection{The unique invariant quadratic form under the bridge}

      By the Rigidity Lemma (Q7 Lemma~4.3), any $\suTwo$-equivariant isomorphism
      $\varphi \colon \Sym^{2}(\Vr) \to \Wsp$ is unique up to a positive scalar $\lambda > 0$.
      The image under $\varphi$ of any $\suTwo$-equivariant endomorphism of $\Sym^{2}(\Vr)$
      is therefore also determined up to $\lambda$.

      \begin{lemma}[Casimir as the unique invariant]\label{lem:casimir-unique}
        The only $\suTwo$-invariant quadratic form on $\Sym^{2}(\Vr)$ is the Casimir
        $Q_{C}(v) = \langle v, \Csu\, v\rangle = 2\lVert v \rVert^{2}$.
        Any other $\suTwo$-equivariant symmetric bilinear form on $\Sym^{2}(\Vr)$ is a scalar
        multiple of $Q_{C}$.
      \end{lemma}

      \begin{proof}
        The module $\Sym^{2}(\Vr)$ is irreducible of dimension~$3$ over $\mathbb{C}$ (the unique
        spin-$1$ representation of $\suTwo$).
        By Schur's lemma, the algebra of $\suTwo$-equivariant endomorphisms
        $\mathrm{End}_{\suTwo}(\Sym^{2}(\Vr)) \simeq \mathbb{C}$:
        the only equivariant endomorphisms are scalar multiples of the identity.
        The Casimir $\Csu = J_{1}^{2} + J_{2}^{2} + J_{3}^{2}$ acts as $j(j+1)\cdot\mathrm{Id} = 2\cdot\mathrm{Id}$
        on the spin-$j = 1$ module, making $Q_{C} = 2\lVert\cdot\rVert^{2}$ the unique
        (up to scalar) invariant quadratic form.
      \end{proof}

    \subsection{The commutator term maps to the Casimir under the bridge}

      \begin{proposition}[Bridge image of the commutator term]\label{prop:bridge-casimir}
        Under the equivariant bridge $\varphi$ of\/ Q7 Conjecture~$4.7$~\cite{Beau2026q7}, the
        sub-principal contribution of\/ $[\tX, \tY] = \tZ$ to the symbol of\/ $\Leff$ is mapped to a
        quadratic form on $\Wsp$ that is proportional to the Casimir $Q_{C}$ restricted to the
        spin-weight-$0$ eigenspace $\{e_{0}\} \subset \Sym^{2}(\Vr)$.
        The proportionality constant is fixed (not a free parameter) by the normalisation
        $\Csu|_{\Sym^{2}(\Vr)} = 2 \cdot \mathrm{Id}$.
      \end{proposition}

      \begin{proof}
        By Proposition~\ref{prop:source}, the sub-principal $k_Z^{2}$-term arises solely from
        $[\tX, \tY] = \tZ$.
        The grading correspondence (Q7 Remark~4.2) identifies $\tZ$ with the spin-weight-$0$
        basis vector $e_{0} \in \Sym^{2}(\Vr)$.
        Under $\varphi$, the $k_Z^{2}$-component of $\sigmatwo(\Leff)$ is carried by
        $\varphi(e_{0}) \in \Wsp$.
        Any $\suTwo$-equivariant quadratic form evaluated on the $e_{0}$-direction must, by
        Lemma~\ref{lem:casimir-unique}, be a scalar multiple $\lambda \cdot 2$ of the Casimir
        eigenvalue on that direction, for some $\lambda > 0$ depending on the normalisation of
        $\varphi$.

        The scalar $\lambda$ is fixed by consistency with the horizontal sector.
        The same equivariant identification $\varphi$ must simultaneously map the
        $e_{\pm}$-directions (spin-weight $\pm 1$, corresponding to horizontal generators $\tX, \tY$)
        to the horizontal sector of $\Wsp$, where the coefficient is $\AH$.
        Since $\Csu$ acts as $2 \cdot \mathrm{Id}$ uniformly on all of $\Sym^{2}(\Vr)$ (the Casimir
        is a central element, so its eigenvalue $j(j+1) = 2$ is the same on all weight subspaces),
        the equivariant form assigns the same scalar $\lambda \cdot 2$ to both the $e_0$ and the
        $e_{\pm}$ directions.
        Therefore $\Az = \lambda \cdot 2 = \AH$, and the ratio $\Az/\AH = 1$ is independent of
        $\lambda$.
        With $\AH \to 2$ proved in Q10~\cite{Beau2026q10} and U1~\cite{Beau2026u1}, the
        normalisation forces $\lambda = 1$ and $\Az = 2$ in the asymptotic regime $q \to \infty$.
        There is no free scalar in the identification $\Az = 2$: $\lambda$ is eliminated by the
        horizontal sector and $\AH$.
      \end{proof}

      \begin{proof}[Proof of Theorem~\ref{thm:main}]
Proposition~\ref{prop:bridge-casimir} gives $\Az = \Csu = 2$.
By Q7 Corollary~6.2 (no cross terms) and Proposition~\ref{prop:bridge-casimir}
($\Az = 2$), the spatial block of the symbol is
\begin{equation}
  \sigmatwo(\Leff)\big|_{\Wsp} = \AH(k_X^{2} + k_Y^{2}) + 2\, k_Z^{2}.
\end{equation}
Q7 Conjecture~6.6~\cite{Beau2026q7} asserted $\AH \to 2$ as $q \to \infty$; this is now
proved in Q10~\cite{Beau2026q10} and U1~\cite{Beau2026u1} under the Q5a hypotheses and the
O-series spectral universality~[U], with rate $|\AH - \Az| = O(q^{-1/2})$.
The full effective co-metric is therefore
\begin{equation}
  g^{\mu\nu} = \mathrm{diag}(-A_{\tau},\, 2,\, 2,\, 2),
\end{equation}
which is non-degenerate of rank~$4$.
Combined with the Lorentzian signature result of Q5b Theorem~6.1~\cite{Beau2026q5b},
this gives signature $(-,+,+,+)$ and closes Q5b-O2 under the Q5a hypotheses and~[U].
      \end{proof}

      \begin{remark}[Unconditional character of $\Az = 2$]
The equality $\Az = 2$ is a consequence of pure representation theory (Schur's lemma applied
to $\Sym^{2}(\Vr)$ and the algebraic identity $\Csu = 2\cdot\mathrm{Id}$ on spin-$1$).
It does not depend on the numerical pipeline, on the choice of prime $q$, or on any
regularity hypothesis beyond bridge non-obstruction.
The Q5a hypotheses [H1], [H-w], [H-E1], [C] are the only foundational conditions:
[H-lift] is a theorem under these (Q9~\cite{Beau2026q9}), and bridge non-obstruction
(exclusion of Case~3) is proved by the coercivity argument of Q9 Corollary~5.2~\cite{Beau2026q9}.
      \end{remark}

  \section{Consistency with Hypoelliptic Symbol Calculus}\label{sec:beals}

    \subsection{The Beals\textendash Greiner framework}

      The theory of hypoelliptic pseudodifferential operators on graded nilpotent groups
      (Beals\textendash Greiner calculus; see also Taylor's operator algebras on nilpotent groups)
      provides an independent setting for computing the sub-principal symbol of $\Leff$ directly.
      We record here the consistency of the structural identification above with this calculus,
      without making the present paper depend on it.

      \begin{proposition}[Beals\textendash Greiner consistency]\label{prop:bg}
        In the Beals\textendash Greiner symbol calculus adapted to $\heisR$, the sub-principal
        symbol of an operator of the form $\Leff = \tX^{2} + \tY^{2} + (\text{lower-order})$
        at homogeneous degree~$2$ in the $\tZ$-direction receives contributions solely from
        the commutator term $[\tX, \tY] = \tZ$.
        No contributions at higher sub-principal order arise within the class-$2$ nilpotency
        of $\heisR$, since the lower central series terminates at step~$2$.
      \end{proposition}

      The proof is a standard computation in the Heisenberg calculus and is recorded for
      completeness; it confirms Proposition~\ref{prop:source} at the level of the symbol
      expansion without adding new input to the identification of $\Az$.

      \begin{remark}[Why Beals\textendash Greiner is not the primary tool]
A direct computation via the hypoelliptic symbol expansion would confirm $\Az = 2$ but
would require significant technical apparatus (Plancherel decomposition over the
Schr\"odinger representation fibres, symbol estimates on $L^{2}(\mathbb{R})$-valued kernels)
that is out of proportion with the simplicity of the structural argument of
Section~\ref{sec:casimir}.
The structural argument is preferred because it makes the reason for $\Az = 2$ transparent:
it is a consequence of Schur's lemma and the Casimir eigenvalue, not an artefact of the
specific form of the expansion.
      \end{remark}

  \section{Consequences for the Effective Metric and Q5b-O2}\label{sec:consequences}

    \subsection{Resolution of Q5b-O2}

      Open problem Q5b-O2 of~\cite{Beau2026q5b} asked for the extension of the effective metric
      to a full $4 \times 4$ Lorentzian tensor by incorporating the sub-principal $\tZ$-correction.
      Theorem~\ref{thm:main} provides this extension:

      \begin{corollary}[Full rank-$4$ effective metric]\label{cor:metric}
        Under the Q5a hypotheses, [U] \textup{(}O-series spectral universality, proved in
        U1~\cite{Beau2026u1}\textup{)}, and the bridge non-obstruction \textup{(}Q7\textendash Q9\textup{)},
        the effective co-metric tensor on $\mathbb{R}_{\tau} \times \heisR$ is
        \begin{equation}
          g^{\mu\nu} = \mathrm{diag}(-A_{\tau},\, 2,\, 2,\, 2),
        \end{equation}
        which is non-degenerate of rank~$4$, with Lorentzian signature $(-,+,+,+)$ and
        with all spatial eigenvalues equal to the $\suTwo$-Casimir eigenvalue $\Csu = 2$.
        The asymptotic isotropy $\AH \to 2$ with rate $O(q^{-1/2})$ is proved in
        Q10~\cite{Beau2026q10} and U1~\cite{Beau2026u1}.
        Open problem $Q5b$-$O2$ is resolved under the Q5a hypotheses and~[U].
      \end{corollary}

    \subsection{The bridge conjecture: status update}

      Q7 Conjecture~4.7~\cite{Beau2026q7} posits the existence of the equivariant bridge
      $\varphi \colon \Sym^{2}(\Vr) \xrightarrow{\sim} \Wsp$.
      Q9~\cite{Beau2026q9} proves the bridge non-obstruction: the pathological Case~3 (obstruction
      to positivity) is excluded by the coercivity of the Mosco limit, leaving only the
      isotropic Case~1 ($\AH = \Az$) and the anisotropic Case~2 ($\AH \neq \Az$).
      Q10~\cite{Beau2026q10} and U1~\cite{Beau2026u1} then prove $\AH \to 2 = \Az$ under the
      Q5a hypotheses and~[U], establishing that only Case~1 survives in the limit.
      The bridge conjecture is thus resolved in the isotropic case conditionally on the Q5a
      hypotheses and~[U].
      Full bridge existence at finite $q$ (Q7 Conjecture~4.7 in its original form) remains an
      open problem at the structural level; what is proved is the unique isotropic limit.

    \subsection{Structural significance: $\Az$ is not a free parameter}

      A central feature of the result is that $\Az$ is not a tunable constant of the effective
      geometry.
      It is fixed by the Casimir of the unique irreducible $\suTwo$-module of dimension~$3$.
      This is consistent with the broader programme principle: physical parameters emerge from
      the spectral structure of the substrate, not from external inputs.
      The effective metric is therefore determined, up to the single temporal coefficient $A_{\tau}$
      (subject to Q5b-O3~\cite{Beau2026q5b}), by the representation-theoretic data already
      established in O23~\cite{Beau2026a23}, O28\textendash O29~\cite{Beau2026a32,Beau2026a33},
      and Q7~\cite{Beau2026q7}.

  \section{Summary and Open Problems}

    \subsection{What this paper establishes}

      \begin{enumerate}
        \item The sub-principal $k_Z^{2}$-term in $\sigmatwo(\Leff)$ has a unique algebraic source:
        the Heisenberg commutator $[\tX, \tY] = \tZ$ (Proposition~\ref{prop:source}).
        \item Under the equivariant bridge, the coefficient of this term is fixed by the
        $\suTwo$-Casimir: $\Az = \Csu = 2$ (Lemma~\ref{lem:casimir-unique},
        Proposition~\ref{prop:bridge-casimir}).
        \item Consequently the effective spatial co-metric is isotropic with all eigenvalues
        equal to~$2$, and the full co-metric $g^{\mu\nu}$ is non-degenerate rank-$4$ Lorentzian
        (Theorem~\ref{thm:main}, Corollary~\ref{cor:metric}).
        \item Open problem Q5b-O2 of~\cite{Beau2026q5b} is resolved under the hypotheses of
        Theorem~\ref{thm:main}.
      \end{enumerate}

    \subsection{Remaining open problems}

      \textbf{Q8-O1} \emph{(Bridge existence at finite $q$).}
      Q9~\cite{Beau2026q9} proves the bridge non-obstruction (Case~3 excluded), and
      Q10\textendash U1 close the isotropic limit $\AH = \Az = 2$.
      What remains open is the finite-$q$ structural question: does an explicit
      $\suTwo$-equivariant isomorphism $\varphi_q \colon \Sym^{2}(\Vr) \xrightarrow{\sim} \Wsp$
      exist at each prime $q$, not just in the asymptotic limit?
      This is now a secondary open problem: the programme consequences are established
      without it.

      \textbf{Q8-O2} \emph{(Unconditional isotropy --- closed).}
      Q7 Conjecture~6.6 (asymptotic $\AH \to 2$ as $q \to \infty$) is now proved in
      Q10~\cite{Beau2026q10} (structural argument under~[U]) and
      U1~\cite{Beau2026u1} (proof of~[U] with rate $O(q^{-1/2})$), both conditionally on the
      Q5a hypotheses.
      This open problem is closed.

      \textbf{Q8-O3} \emph{(Identification of $A_{\tau}$ and completion of Q5b-O3).}
      With $\Az = 2$ established, the remaining free coefficient is $A_{\tau}$
      (the temporal sector of the effective metric).
      Its identification from the Born\textendash Infeld saturation data is the subject of
      Q5b-O3~\cite{Beau2026q5b}.

      \clearpage
      \bibliography{cosmochrony-bibliography}

\end{document}
